\documentclass[12pt]{article}
\usepackage{geometry}                % See geometry.pdf to learn the layout options. There are lots.
\geometry{letterpaper}                   % ... or a4paper or a5paper or ... 

\usepackage[parfill]{parskip}    % Activate to begin paragraphs with an empty line rather than an indent
\usepackage{fullpage}
\usepackage{graphicx}
\usepackage{grffile}
\usepackage{listings}
\usepackage{hyperref}
\usepackage{tabularx} %tabular with stretch columns
\usepackage{enumerate}
\usepackage{pdfpages}
\usepackage{pdflscape}
\usepackage[utf8]{inputenc}

\usepackage[acronym]{glossaries} % make a separate list of acronyms
\makeglossaries

%testing
% Index
\usepackage{makeidx}
\makeindex

\begin{document}


\nocite{Sh:1}

\begin{titlepage}

\begin{minipage}{4cm}
\begin{tabular}{l}
\includegraphics[width=0.5\textwidth]{Images/uppsala}
\end{tabular}
\end{minipage}
\hfill
\begin{minipage}{4cm}
\begin{tabular}{r}
\includegraphics[width=0.5\textwidth]{Images/logo}
\end{tabular}
\end{minipage}


\begin{center}
% Upper part of the page

\textmd{HealthShare Final Report}
\vfill
% Title
 \huge{\textbf{A Study on Methods to Increase Interoperabilty and Unify Electronic Healthcare Records } }\\[2.0cm]
for\\
\large\textbf{Uppsala County Council}\\[1.0cm]
by\\
\large{\textbf{Uppsala Universitet}} and \large{\textbf{Rose-Hulman Institute of Technology}}\\[1.0cm]

\end{center}

\begin{center}
\vfill
December 2011\\
\end{center}

\end{titlepage}

\tableofcontents
\newpage

\begin{abstract}
Our abstract goes here... 
\end{abstract}

\newpage

\printglossaries

\newpage

\section{Introduction}
Define the problem and in what context (maybe contextualize, write a story, describe a scenario? a concrete relevant scenario that occurs with persons involved)


\subsection{Background}

Romel hah The setting of problems in EHR/EMR interoperability situations. Describe what the problems related to interoperability are within healthcare. The situation in Uppsala County, why we are doing this. What is being done on interoperability in other areas. Describe the problems related to interoperating EHR.
describe the ideal interoperating system

\subsection{Purpose and Scope}
To investigate the interoperability of
EHR/EMR systems taking in regard what
the people involved want, the limitations,
standards and organizational issues.

Explain what we will look into and why, limit the scope to what we will be discussing in the paper.

\subsection{Reading Instructions}
instructions for different audiences, so people who already know, or don't care about certain sections, don't have to sift through the paper to find information relevant to them. For the rest of the report.

\newpage

\section{Method}

\subsection{Project Organization}
 (How, and our own organization)

\subsection{Interviews}
(preparations, structure, outlines/protocols, focus of the interviews)
\begin{itemize}
\item phone interviews
\item regular interviews
\item hybrids between the two
\end{itemize}

How did we prepare the interviews, how they where/wasn't structured, transcibed(?).

\subsection{Information search}
Reading, methods for searching for information

\newpage

\section{What the people involved want}
(maybe move to introduction if too small)
\begin{itemize}
\item physicians
\item nurses
\item patients
\item administrative personal
\item the public
\item politicians
\item technicians
\item legal authorities
\item lab workers
\end{itemize}

\newpage

\section{Interoperabilty}

\subsection{Introduction} % Introduction to what this chapter should be about. (relevant, not too long, not too short) ...Should this really be a subsection or just text?

\subsection{Definition within healthcare} %Define what interoperability is. Talk about interoperability in general, and also in this specific case. Why interoperability is important and how it works.

\subsection{Current situation} %What the current situation is regarding interoperability. Why it is a big problem.
\subsection{Prerequisites for interoperability on a national level}

To achieve European wide solutions for cross-border transfer of patient information certain requirements need to be fulfilled. There has to be an agreement on the definitions of data sets for both patient summaries and e-prescription, a legal framework for data transfer, a technical framework to connect the systems at each level and a working semantic interoperability. \cite{epSOS1}

\subsection{Standardized care plans}

%Subsections on different interoperability solutions.
\subsection{Solution 1}
\subsection{Solution 2}
\subsection{Solution 3}

\newpage

\section{Technical Standards}


\subsection{What is a Technical Standard for Computer Software?}
A technical standard is a recognized and established requirement about software systems that establishes the ``way things should be done." They are responsible for everything from the uniformity of web browsers (although not all browsers strictly conform to the standard) to the ability for nearly any laptop to connect wirelessly to any access point. Such standards specify aspects of a program such as a commont format for a file or data transfer that allow different developers to develop separate software programs while allowing for interoperability between the different systems. Software standards are the foundation for the interoperability of different software systems.

\subsection{The Creation of a Standard}
Typically for different parties (software development companies) to agree on a specific software standard they create a software standards organization that consist of members and representatives of the various software companies who contribute ideas and opinions about making a single, unifed standard addressing the data interoperability problem that needs to be addressed.

\subsubsection{ Examples of Software Standards Organizations}

\begin{itemize}
\item World Wide Web Consortium (W3C) - Responsible for web standards such as HTML, HTTP, and XML 
\item Institute of Electrical and Electronics Engineers Standards Association (IEEE Standards Association) - Responsible for a wide range of standards for engineering such as the 802.11 standard
\item Internet Engineering Task Force (IETF) - Responsible for providing standards for new internet related technology 
\end{itemize}

\subsection{ Adherance to Standards}
Adherance to a particular software standard can be either manditorally required or volunteerly followed. For example, in the United States banking software must conform to security standards and regulations set forth by the Federal Deposit Insurance Corporation (FDIC). Any system that does not conform to these standards faces legal action by the federal government and as such is not allowed for use by U.S. banks. On the other end of the spectrum, the standards set forth by the W3C as to how a web browser should interperet and render various webpages is a completely voluntary standard and while most popular, modern browsers do conform in the most part to these standards there are many cases where some (if not all) do not completely achieve complience. However, they face no legal retaliation for not complying with the W3C's standards as the standards are volunteerly adhered to and not legally mandated.


\subsection{Software Standards and EHR Software}
With the growing ubiquitousness of the internet, computers, and modern technology as well as the geographic diversity of medical tallent and information it is doubtless that someday a standard specifying the transmission, storage, and format of electronic healthcare records will be developed; the benefit of such a standard is too great to be ignored for long. Such a standard will require and large organization to oversee and develop as patent data is widely varried and medical breakthroughs are discovered daily requiring such a standard to be constantly updated and modified to stay relevent. 

\subsection{The Alternative to a Future Standard For EHR Software}
The only other option for sharing healthcare records would be to have a unified software system doctors all over the world would use to view and modify patient records. However, such a solution is unrealistic as for one there are already too many EHR software companies to resonably plan for one, new company to globally take over the entire electronic healthcare record industry. Furthermore, with anti-trust laws as well as fair competition clauses in many countries (such as the United States) a single organization would not be legally allowed to be the sole producer of EHR software systems.

\newpage

\section{Adoption of future systems}

%Current limitations and obstacles

\subsection{Technical}
Security? Limitations for what that can technically be done?


\subsection{Legal}
What does the laws allow? 

The new Swedish legislation provides new opportunities, but not obligations. \cite{RiR19}

\subsection{Organizational}

== current organizational issues that limit the adoption of interoperating systems. ==

%Here are related items from the Empirica interview
\begin{itemize}
\item When it comes to interoperability, technical issues and standard specifications are subordinated to legal and organizational issues. The organizational issues include common understanding from people involved in the process and getting the commitment from doctors and nurses. 
\item Clear objectives and strong project leaders with the backing of the top level management and politics are very important factors of success for EHR projects. Continuous development and the ability to revise things are also crucial.
\item It is important to reach all the layers of the health care organization to achieve change, to engage the different groups of doctors, nurses, secretaries, technicians and not only have the software developers talk to the management section of the organization.
\item Electronic systems are changing the way doctors work and what responsibilities they have.
\end{itemize}

\newpage

\section{Results/ Discussion}

\subsection{ Prerequisites for increased interoperability}

\subsection{Types of solutions and implications for the end user}

\newpage

\section{Conclusions}

\section{How to do Glossary Entries}

Look at the {\LaTeX} for this section to see how add glossary entries and how to use the in the text.

\newglossaryentry{sample}{name={sample}, description={a sample entry}}

\newacronym[\glsshortpluralkey=cas,\glslongpluralkey=contrived acronyms]{aca}{aca}{a contrived acronym}

A \gls{sample} entry and \gls{aca}. Second use: \gls{aca}.

Plurals: \glspl{sample}. Reset acronym\glsreset{aca}.
First use: \glspl{aca}. Second use: \glspl{aca}.



\newpage

\begin{appendix}


\end{appendix}

\newpage

\bibliographystyle{unsrt} 
\bibliography{ourBib} 

\end{document} 
